\documentclass[10pt,a4paper]{amsart}
\usepackage{amsmath,amssymb,amsthm,graphicx}
\usepackage{verbatim}
\newenvironment{reference}{\comment}{\endcomment}
%\newenvironment{reference}{}{}
\newenvironment{slogan}{\comment}{\endcomment}
\newenvironment{history}{\comment}{\endcomment}
% For commutative diagrams we use Xy-pic
\usepackage[all]{xy}

% We use 2cell for 2-commutative diagrams.
\xyoption{2cell}
% Theorem environments.
%
\theoremstyle{plain}
\newtheorem{theorem}[subsection]{Theorem}
\newtheorem{proposition}[subsection]{Proposition}
\newtheorem{lemma}[subsection]{Lemma}

\theoremstyle{definition}
\newtheorem{definition}[subsection]{Definition}
\newtheorem{example}[subsection]{Example}
\newtheorem{exercise}[subsection]{Exercise}
\newtheorem{situation}[subsection]{Situation}

\theoremstyle{remark}
\newtheorem{remark}[subsection]{Remark}
\newtheorem{remarks}[subsection]{Remarks}

\numberwithin{equation}{subsection}

% Macros
%
\def\lim{\mathop{\mathrm{lim}}\nolimits}
\def\colim{\mathop{\mathrm{colim}}\nolimits}
\def\Spec{\mathop{\mathrm{Spec}}}
\def\Hom{\mathop{\mathrm{Hom}}\nolimits}
\def\Ext{\mathop{\mathrm{Ext}}\nolimits}
\def\SheafHom{\mathop{\mathcal{H}\!\mathit{om}}\nolimits}
\def\SheafExt{\mathop{\mathcal{E}\!\mathit{xt}}\nolimits}
\def\Sch{\mathit{Sch}}
\def\Mor{\mathop{\mathrm{Mor}}\nolimits}
\def\Ob{\mathop{\mathrm{Ob}}\nolimits}
\def\Sh{\mathop{\mathit{Sh}}\nolimits}
\def\NL{\mathop{N\!L}\nolimits}
\def\CH{\mathop{\mathrm{CH}}\nolimits}
\def\proetale{{pro\text{-}\acute{e}tale}}
\def\etale{{\acute{e}tale}}
\def\QCoh{\mathit{QCoh}}
\def\Ker{\mathop{\mathrm{Ker}}}
\def\Im{\mathop{\mathrm{Im}}}
\def\Coker{\mathop{\mathrm{Coker}}}
\def\Coim{\mathop{\mathrm{Coim}}}

% Boxtimes
%
\DeclareMathSymbol{\boxtimes}{\mathbin}{AMSa}{"02}

%
% Macros for moduli stacks/spaces
%
\def\QCohstack{\mathcal{QC}\!\mathit{oh}}
\def\Cohstack{\mathcal{C}\!\mathit{oh}}
\def\Spacesstack{\mathcal{S}\!\mathit{paces}}
\def\Quotfunctor{\mathrm{Quot}}
\def\Hilbfunctor{\mathrm{Hilb}}
\def\Curvesstack{\mathcal{C}\!\mathit{urves}}
\def\Polarizedstack{\mathcal{P}\!\mathit{olarized}}
\def\Complexesstack{\mathcal{C}\!\mathit{omplexes}}
% \Pic is the operator that assigns to X its picard group, usage \Pic(X)
% \Picardstack_{X/B} denotes the Picard stack of X over B
% \Picardfunctor_{X/B} denotes the Picard functor of X over B
\def\Pic{\mathop{\mathrm{Pic}}\nolimits}
\def\Picardstack{\mathcal{P}\!\mathit{ic}}
\def\Picardfunctor{\mathrm{Pic}}
\def\Deformationcategory{\mathcal{D}\!\mathit{ef}}
\usepackage{cleveref}
\crefname{originallemma}{original lemma}{original lemmas}
\Crefname{originallemma}{Original lemma}{Original lemmas}
\makeatletter
\newlabel{divisors-lemma-c1-additive}{{02SL}{source}{}{sourcepointer-1}{}}
\newlabel{divisors-lemma-c1-additive@cref}{{[originallemma][0][]02SL}{[1][0][]source}{}{}{}}
\makeatother
\crefname{originalequation}{original equation}{original equations}
\Crefname{originalequation}{Original equation}{Original equations}
\makeatletter
\newlabel{equation-equal-as-cycles}{{02SR}{source}{}{sourcepointer-2}{}}
\newlabel{equation-equal-as-cycles@cref}{{[originalequation][0][]02SR}{[1][0][]source}{}{}{}}
\makeatother
\crefname{originallemma}{original lemma}{original lemmas}
\Crefname{originallemma}{Original lemma}{Original lemmas}
\makeatletter
\newlabel{lemma-compute-c1}{{02SK}{source}{}{sourcepointer-3}{}}
\newlabel{lemma-compute-c1@cref}{{[originallemma][0][]02SK}{[1][0][]source}{}{}{}}
\makeatother
\crefname{originallemma}{original lemma}{original lemmas}
\Crefname{originallemma}{Original lemma}{Original lemmas}
\makeatletter
\newlabel{lemma-pullback-coherent}{{02RE}{source}{}{sourcepointer-4}{}}
\newlabel{lemma-pullback-coherent@cref}{{[originallemma][0][]02RE}{[1][0][]source}{}{}{}}
\makeatother
\crefname{originalreference}{original reference}{original references}
\Crefname{originalreference}{Original reference}{Original references}
\makeatletter
\newlabel{situation-setup}{{02QL}{source}{}{sourcepointer-5}{}}
\newlabel{situation-setup@cref}{{[originalreference][0][]02QL}{[1][0][]source}{}{}{}}
\makeatother
\title{Flat pullback of the divisor of a meromorphic section of an invertible sheaf}
\author{The Stacks Project Authors}
\date{stacks-project:chow.tex:local-snapshot-20261012; GNU FDL 1.2 or later}
\begin{document}
\maketitle
\noindent\textit{Editorial note.} All delivered mathematical text is the author's own, copied from the pinned Stacks Project chapter file chow.tex; no statement or proof step was rewritten, reordered or removed.\par\medskip
\noindent\textit{Editorial note.} Preamble provenance: the pinned chapter file opens with \textbackslash{}input\{preamble\} and carries no preamble text of its own, so the project's single Stacks-wide preamble staged at sources/stacks-chow/preamble.tex is excerpted verbatim as three spans, lines 8-12 (the verbatim package together with the project's own reference, slogan and history display-environment definitions, including the author's commented-out older form of the reference environment), lines 14-18 (Xy-pic setup) and lines 154-217 (the native theorem declarations and mathematical macros). The lines 8-12 span is required because a counted statement or its surrounding text may use those display environments, which the project typesets as comments. Nothing was added, removed or redefined; the wrapper class and the amsmath/amssymb/amsthm/graphicx base packages are supplied by the assembler.\par\medskip
\noindent\textit{Editorial note.} Same-chapter pointers are located at their \textbackslash{}label lines inside the pinned chapter file; cross-chapter pointers are located in the project's own tag list staged in the snapshot (sources/stacks-chow/tags.txt), at the verbatim lines carrying their canonical tags.\par\medskip
\noindent\textit{Allowed standard prerequisites (editorial).}\begin{itemize}
\item Additivity of the divisor of a meromorphic function. Divisors, label divisors-lemma-c1-additive, invoked at proof line 4514. Reduces the statement for a general meromorphic section written as a quotient of two global sections to the two statements for those sections.
\item The cycle identity for the zero scheme of a section and its restriction to the components. Same chapter chow.tex, the equation labelled equation-equal-as-cycles, invoked at proof line 4517 and established by S14-algebra-prepare-geometric-cap. Identifies the class of the zero scheme of the pulled back section with the sum over the components of the divisors of the restricted sections.
\item The cycle class of a divisor and flat pullback of cycle classes of coherent sheaves. Same chapter chow.tex, labels lemma-compute-c1 (proof line 4528) and lemma-pullback-coherent (proof line 4528, claimed as S14-algebra-pullback-coherent). Rewrites the flat pullback of the divisor as the flat pullback of a cycle class and then as the class of the inverse image, which coincides with the zero scheme of the pulled back section.
\end{itemize}\medskip
\ifdefined\hypertarget\hypertarget{sourcepointer-1}{}\fi
\noindent\textit{Source reference.} Original lemma 02SL: sources/stacks-chow/tags.txt, source lines 3664--3664. Navigation pointer for \textbackslash{}ref\{divisors-lemma-c1-additive\} used inside the counted span (line 4514); the argument that consumes it is retained in the span.\par\medskip
\ifdefined\hypertarget\hypertarget{sourcepointer-2}{}\fi
\noindent\textit{Source reference.} Original equation 02SR: sources/stacks-chow/pinned-original-source.tex, source lines 4375--4375. Navigation pointer for \textbackslash{}ref\{equation-equal-as-cycles\} used inside the counted span (line 4517); the argument that consumes it is retained in the span.\par\medskip
\ifdefined\hypertarget\hypertarget{sourcepointer-3}{}\fi
\noindent\textit{Source reference.} Original lemma 02SK: sources/stacks-chow/pinned-original-source.tex, source lines 4198--4198. Navigation pointer for \textbackslash{}ref\{lemma-compute-c1\} used inside the counted span (line 4528); the argument that consumes it is retained in the span.\par\medskip
\ifdefined\hypertarget\hypertarget{sourcepointer-4}{}\fi
\noindent\textit{Source reference.} Original lemma 02RE: sources/stacks-chow/pinned-original-source.tex, source lines 2331--2331. Navigation pointer for \textbackslash{}ref\{lemma-pullback-coherent\} used inside the counted span (line 4528); the argument that consumes it is retained in the span.\par\medskip
\ifdefined\hypertarget\hypertarget{sourcepointer-5}{}\fi
\noindent\textit{Source reference.} Original reference 02QL: sources/stacks-chow/pinned-original-source.tex, source lines 1372--1372. Navigation pointer for \textbackslash{}ref\{situation-setup\} used inside the counted span (line 4485); the argument that consumes it is retained in the span.\par\medskip
\begin{lemma}
\label{lemma-prepare-flat-pullback-cap-c1}
Let $(S, \delta)$ be as in Situation \ref{situation-setup}.
Let $X$, $Y$ be locally of finite type over $S$.
Let $f : X \to Y$ be a flat morphism of relative dimension $r$.
Let $\mathcal{L}$ be an invertible sheaf on $Y$.
Assume $Y$ is integral and $n = \dim_\delta(Y)$.
Let $s$ be a nonzero meromorphic section of $\mathcal{L}$.
Then we have
$$
f^*\text{div}_\mathcal{L}(s) = \sum n_i\text{div}_{f^*\mathcal{L}|_{X_i}}(s_i)
$$
in $Z_{n + r - 1}(X)$. Here the sum is over the irreducible
components $X_i \subset X$ of $\delta$-dimension $n + r$,
the section $s_i = f|_{X_i}^*(s)$ is the pullback of $s$, and
$n_i = m_{X_i, X}$ is the multiplicity of $X_i$ in $X$.
\end{lemma}

\begin{proof}
To prove this equality of cycles, we may work locally on $Y$.
Hence we may assume $Y$ is affine and $s = p/q$ for some nonzero
sections $p \in \Gamma(Y, \mathcal{L})$ and $q \in \Gamma(Y, \mathcal{O})$.
If we can show both
$$
f^*\text{div}_\mathcal{L}(p) =
\sum n_i\text{div}_{f^*\mathcal{L}|_{X_i}}(p_i)
\quad\text{and}\quad
f^*\text{div}_\mathcal{O}(q) =
\sum n_i\text{div}_{\mathcal{O}_{X_i}}(q_i)
$$
(with obvious notations) then we win by the
additivity, see Divisors, Lemma \ref{divisors-lemma-c1-additive}.
Thus we may assume that $s \in \Gamma(Y, \mathcal{L})$.
In this case we may apply the equality
(\ref{equation-equal-as-cycles}) to see that
$$
[Z(f^*(s))]_{k + r - 1} =
\sum n_i\text{div}_{f^*\mathcal{L}|_{X_i}}(s_i)
$$
where $f^*(s) \in f^*\mathcal{L}$ denotes the pullback of $s$ to $X$.
On the other hand we have
$$
f^*\text{div}_\mathcal{L}(s) = f^*[Z(s)]_{k - 1}
= [f^{-1}(Z(s))]_{k + r - 1},
$$
by Lemmas \ref{lemma-compute-c1} and \ref{lemma-pullback-coherent}.
Since $Z(f^*(s)) = f^{-1}(Z(s))$ we win.
\end{proof}

\end{document}
